% ===========================================================
\section{Discussion}
% ===========================================================

\subsection{Performance Analysis}

\TODO{Discuss which chemical families (e.g., Cr-based, Zr-based)
show the highest/lowest error and why.}

The largest errors are observed for MAX phases containing
heavy A-group elements (e.g., Bi, Tl) whose soft acoustic
branches are sensitive to small changes in the interatomic
potential.

\subsection{Role of Elastic Constants}

Including elastic constants as global graph features reduces
MAE by \TODO{X\%}, confirming that macroscopic mechanical
properties encode information complementary to the local
atomic environment.

\subsection{Comparison with eSEN Methodology}

While eSEN learns the full energy surface $E(\mathbf{r},\mathbf{a})$
and derives IFCs analytically via $\nabla^2 E$,
our approach directly predicts IFCs from graph embeddings.
This trade-off offers faster training at the cost of losing
the strict energy-conservation guarantee of force-field models.
Future work will explore hybrid approaches.

\subsection{Limitations}

\begin{itemize}
  \item The dataset is limited to M\textsubscript{2}AX stoichiometry;
        extension to M\textsubscript{2}AX\textsubscript{2} and other
        variants is left for future work.
  \item Imaginary (unstable) modes are predicted as negative
        frequencies; explicit modelling of soft-mode instabilities
        warrants further investigation.
  \item The model is trained on DFT-LDA/GGA data and may
        not transfer to hybrid-functional or beyond-DFT results.
\end{itemize}
